---
format: yohaku/0.2
target: text
macro: '[[macros/raw]]'
---

## 正文
````text
# Yohaku

给人与 AI 的本地工作空间。
A local workspace for people and AI.

## ⬇️ [下载 Windows 版 / Download for Windows](https://github.com/epsilon-lain/yohaku/releases/download/v2026.11.17-preview1/Yohaku-2026.11.17-preview1-windows-x64.zip)

**Windows 11 x64 · 约 15 MB · 便携技术预览 2026.11.17-preview1**
无需另装 Rust、Python 或 WSL。No Rust, Python or WSL installation required.

1. **下载 / Download**：点击上方链接，保存 Windows ZIP。
2. **解压 / Extract**：完整解压到有写入权限的文件夹，保留所有包内文件。
3. **启动 / Launch**：双击解压目录中的 **`Launch-Yohaku.cmd`**。

[首次使用与常见问题 / Getting started](docs/windows/README.md) · [版本说明与校验 / Release notes](https://github.com/epsilon-lain/yohaku/releases/tag/v2026.11.17-preview1) · [反馈问题 / Report an issue](https://github.com/epsilon-lain/yohaku/issues)

第一次打开可以先体验本地文档、Loom 和 Socius 设置。模型连接可稍后自行配置；
真实模型工具调用和图片理解尚未验收。资料默认保存在解压目录的 `data/Yohaku`。

请使用上方 Windows 下载链接；GitHub 的 **Code → Download ZIP** 和 **Source code** 是源码归档。
**仅在当前电脑隔离环境验证，尚未在独立干净 Windows 环境验证。**

---

## 开发与研究 / Development

以下构建与能力说明对应 Linux / WSL 源码。Windows 预览的实际范围、运行包和独立源码身份见
[Windows 使用说明](docs/windows/README.md)及 Release 附件 `BUILD-SOURCE.json`。

Yohaku is an experimental native, local-first workspace for people and AI.
It is independent of Lain Brain. Loom, the deterministic `.tex` text generator,
is kept in this repository. This source snapshot is **not open source**:
original material is all rights reserved. See [LICENSE](LICENSE) and
[THIRD_PARTY_NOTICES.md](THIRD_PARTY_NOTICES.md).

## Capabilities

- Rust/Smithay native documents, local tools and Socius chat.
- Loom-managed body edits that preserve notes, macros and output routing.
- Built-in entry listing, name/description search, stable node IDs, bounded BFS
  and a rebuildable `.tex`-derived index; no plugin required.
- Versioned Unicode line/column reads and guarded coordinate edits.
- Registered tools with actual exit codes/logs in an offline bubblewrap sandbox.
- Existing optional model connectors, guarded independent Codex tasks and model selection.
- Native terminal sessions/windows, scoped clipboard transfer and nested text input.
- codex-text/3 content/send/stop pages with save-feedback/2 operation-scoped receipts.
- A fixed internal candidate prepare/check/activate/status/rollback workflow.

## Requirements

Tested on Arch Linux in WSL2, Windows 11, with the X11 host backend. Use an
existing working graphics/WSLg setup; this project does not install or repair it.
The separate native Windows technical preview has only the verification scope stated above; macOS builds are not verified.

Requirements: Rust/Cargo supporting edition 2024, Python 3, C compiler, binutils,
pkg-config, Wayland client/server libraries, EGL/GLES2, libxkbcommon, X11
libraries and a usable installed font. Bubblewrap is required to run world
tools; there is no unsandboxed fallback. Exact tested versions and results are
in [VERIFICATION.md](VERIFICATION.md).

## Build and run

From this repository in WSL/Linux:

```sh
./scripts/build.sh
python3 scripts/init_config.py
./run-yohaku
```

The build script compiles Loom from `loom-src` with its bundled dependencies,
then the compositor, workspace CLI, terminal window and sample tool window. The
chat-workspace name is a second local copy of the same identified workspace CLI. Locally built
executables go into ignored `bin/` paths. No precompiled executable is committed.
Other Yohaku dependencies use `Cargo.lock` and the normal Cargo cache; fetching
them may require network access on a new machine. With an existing complete
cache, use `CARGO_NET_OFFLINE=true ./scripts/build.sh`.

The launcher selects X11 on WSL, forwards Fcitx variables and requires `DISPLAY`.
`YOHAKU_DISPLAY_BACKEND=native ./run-yohaku` selects the existing native backend.
Always start through the wrapper: runtime paths are relative to the project.

The initializer copies English templates only if local destinations are absent.
It imports no account, login, world or model configuration. Building and starting
send no model request. Configure your own endpoint/model in Socius when needed;
provide your own key through the configured environment variable or session UI.
Never commit it. The experimental local Codex connector requires a separately
installed compatible runtime and explicit human setup/login. Clean-config
verification does not authenticate or exercise a provider.

The text service and trusted host updater ship disabled in their portable configs.
Their blank templates contain no request history, login or machine-specific hashes.
For first-time opt-in setup, permissions, candidate preparation and rollback, see
[docs/INTERNAL-UPDATE.md](docs/INTERNAL-UPDATE.md). Building or installing draft
pages does not send a model request. Each text send is an independent model task;
ordinary Socius chat remains a separate route. A saved revision, an accepted
submission and a completed execution are separate facts. Unknown outcomes must
be queried before another explicit operation.

The current source snapshot includes the actually loaded codex-text/3 and
save-feedback/2 Python implementation. Historical real-model acceptance on the
original installation is recorded separately from this checkout's offline
build/tests in VERIFICATION.md. No account or authenticated session is supplied.

## Minimal node demo

```sh
python3 scripts/create_demo.py --root worlds/demo
python3 agent/yohaku_agent.py --root worlds/demo --config config/ai.json tool list_entry_nodes '{}'
python3 agent/yohaku_agent.py --root worlds/demo --config config/ai.json tool walk_nodes '{"id":"demo:entry","max_depth":3,"limit":10}'
python3 worlds/demo/files/node-demo/test_greet.py
YOHAKU_WORKSPACE="$PWD/worlds/demo" ./run-yohaku
```

Three portable `.tex` documents are imported using existing write,
adopt-existing-source and Loom sync operations. Existing demo targets/tool IDs
are refused, never overwritten. One fixed demo test tool is registered.
No real Lain Brain directory is read or required.

In Chat, ask: “From the Lain Brain demo entry, follow node links to the greeting
module and test. Use coordinate reads to verify the greeting. Do not search
directories instead of using node navigation.” Explicitly authorize any editing
exercise and test run, then inspect revisions, Loom results, exit codes and logs.
Codex chat names have a `yohaku_` prefix; internal/CLI names do not.
Enable the existing Write/Run permissions in Socius before an editing/test
exercise; wording in a prompt does not override disabled tool permissions.

## Source documents and tests

`source-documents.json` lists sibling authoritative `.tex` sources. Generated
text is checked in for bootstrapping. Edit those documents using the existing
Loom format and raw macro, then run:

```sh
python3 scripts/sync_sources.py
python3 scripts/sync_sources.py --check
./scripts/test.sh
```

`--check` generates into a temporary root and compares bytes without replacing
tracked outputs. Demo documents are world inputs, not compiler inputs. Typed
node links never add implicit Loom dependencies. See `loom-src/FORMAT.md` and
`loom-src/tests`. Vendor trees retain upstream formats and notices.

## Known limits

- This prototype is not a security audit or a stable cross-platform release. The separate Windows binary is a limited technical preview.
- Node queries scan `.tex` metadata. BFS has depth/result caps, deduplicates
  cycles and reports broken/ambiguous links; there is no vector database.
- Coordinates are 1-based Unicode characters and end-exclusive. Explicit
  endpoints outside a document are errors; omit them for pagination. Re-read
  after revision conflicts. File-edit success is not program-test success.
- Runtime worlds, logs and actual configuration are private local state.
  Chat summaries do not persist every structured tool argument.
- Physical IME/font/close-button and real AI acceptance were performed on the
  original installation. The clean-config smoke test is separate evidence.
- Font fallback uses fixed installed paths, including Windows FangSong,
  Segoe UI Symbol and Microsoft YaHei in WSL. Linux-only CJK/emoji coverage
  is not guaranteed. Fcitx candidate fonts are configured separately in Fcitx
  and fontconfig. No font is bundled or licensed by this project.
- The existing X11 winit refocus patch is pinned in `vendor/winit` for
  reproducibility, not claimed as the cause of the IME repair. The earlier
  input failure involved a keyboard layout rather than the real Pinyin engine.
  No global Cargo-cache modification is required by this checkout.
- Tool launches use fixed registered arguments. Unfinished parameter changes
  in the private development world remain untouched and are not advertised as
  released functionality. Develop support remains experimental.
- The initial snapshot's two application-discovery fixture failures are preserved
  as historical results in VERIFICATION.md. The refreshed fixtures isolate native
  discovery and assert the current ELF rejection; current results are reported
  separately. No real model, clipboard desktop or physical IME test is implied.
- `apply_update.py` is retained for compatibility tests, not as an upgrade
  payload for arbitrary installations. Build this checkout independently;
  do not use it to replace another running directory.

See [PROVENANCE.md](PROVENANCE.md), [CHANGELOG.md](CHANGELOG.md) and
[VERIFICATION.md](VERIFICATION.md) for scope and measured results.
## Research workflow materials

The separately scoped [workflow runtime snapshot and reproduction instructions](research/workflow-runtime/README.md) include the historical Sol controller, an Adaption candidate source export, host budget and recovery boundaries, and offline checks. They are not enabled in the default desktop product. A [redacted review evidence index](research/review-evidence/README.md) and the [accepted handoff delivery](research/workflow-delivery/README.md) are available for the competition manuscript. No private runtime world, account, credentials or frozen acceptance answers are supplied.

```sh
python3 -B research/workflow-delivery/verify_offline.py
python3 -B research/review-evidence/verify.py
python3 -B -m unittest discover -s research/workflow-runtime/tests -v
```

## π research workflow observations (2026-11-12)

The [bounded π research case note](research/pi-workflow-2026-11-12/README.md) now includes the completed v05 trial: the original owner selected and ran its previously saved verifier, read the result, saved a standalone note, and ended after 12 of 16 allowed requests. Historical memory and external review shaped this continuation; no helpers were created. Independent review confirms a finite numerical witness but still finds a formula error in the revised response. Stable differentiation, newly invented workflows, and causal benefits remain unestablished. This archived case is separate from the earlier multi-agent review trials and does not activate the isolated runtime repairs in the default product.

## Immune-inspired and autonomous workflows (2026-11-16)

The [protocol and first archived observations](research/immune-autonomous-2026-11-16/README.md) compare an immune-inspired organization with a self-organizing lead on the same mathematical seed. Mathematical knowledge is shared immediately while organizational methods remain group-specific, making this a coupled descriptive study. The 18:45:59 China Standard Time snapshot records a saved two-route research plan and a requested continuation in A, and a successful source read after tool-contract friction in B. It establishes neither stable differentiation nor a new mathematical result. The note includes a provenance-linked public evidence projection and links to the separate mathematical formalization and exact graph verification.

## Concept compression and an AC pilot (2026-11-16)

The [first concept-compression pilot](research/concept-compression-ac-2026-11-16/README.md) connects the user's mathematical framework to costed, expandable search macros. A bounded comparison records successes and budget exhaustion for primitive and macro searches, with no shorter solution found by the macro arm in this sample. Four scored AC/Stable AC paths pass the pinned official local verifier; a separate exhaustive certificate proves that ac-08491 has global AC distance eight. These paths are prepared for submission and have not been uploaded to SAIR. The directory includes reproducible code, exact results, verification receipts and an explicit [human/AI collaboration record](research/concept-compression-ac-2026-11-16/COLLABORATION.md). The [separate mathematical note](https://github.com/epsilon-lain/lain-s-note/tree/main/research/concept-compression-2026-11-16) proves finite expression-cost results and the macro-expansion interface. This pilot is separate from the earlier immune/autonomous experiment.

## Two research teams on Sol High/Max (2026-11-16)

The [theory/application research stage](research/theory-application-sol-2026-11-16/README.md) records two explicitly configured Sol/Max leads and a Sol/High external reviewer. The bounded first sprint is complete and checkpointed: the theory team proves scoped positive-property identification results; the application team improves ac-07445 from 17 to 13 primitive moves and certifies its exact ordinary AC distance. The [workflow observation](research/theory-application-sol-2026-11-16/observation/REPORT.md) distinguishes method changes, external interventions and versioned reads; both teams remain single leads, with no stable role differentiation or adoption of the other team's new theorem established. Shared artifacts, snapshots, continuation checkpoints and a locally verified four-line submission are included. This is a Work collaboration run; original Yohaku host migration has not been performed. The [host handoff](research/theory-application-sol-2026-11-16/HOST_HANDOFF.md) preserves the Sol/High floor, Sol Max/Ultra ceiling and prohibition on Astra or weaker fallback.

## Native Yohaku research-team launch handoff (2026-11-16)

The [native launch preparation](research/native-two-team-launch-2026-11-16/README.md) preserves all prior results and provides the [local Codex execution task](research/native-two-team-launch-2026-11-16/CODEX_TASK.md), a 12-item versioned mathematical seed, source verification and an archived-interface audit. The task requires an actual persistent Yohaku research lead to form and dispatch the theory/application directions with Sol/High or above, Sol/Max leads and no Astra or weak fallback. It protects prior in-flight requests and the sealed zero-inference send block. This package is prepared and verified only; no original-host process, native import or research-model request was performed here. Local implementation and real native receipts are still required.

## Native Yohaku theory/application results (2026-11-16)

The [native-run audit and results](research/native-theory-application-2026-11-16/README.md) review knowledge prefix 465 and gate prefix 220 from the user-supplied host archive. A persistent native lead actually created and dispatched the theory and application leads; all three initial tasks completed with requested and acknowledged `gpt-6.1-sol/max` configuration. The native theory derives irredundant-certificate criteria and a family whose deletion cost approaches twice the optimum in a scoped graph model. The application refines the known eight-move ac-08491 witness from work 49 to 48, without shortening it. Later external review reproduces the bounded results and validates the new path with pinned official local code and CLI; no platform upload is claimed. [Workflow evidence](research/native-theory-application-2026-11-16/WORKFLOW.md) establishes an actual read/review/adoption chain and a communication-driven cancellation of planned duplicate work. Three nodes and a single prefix do not establish stable spontaneous differentiation or token savings. The root continuation was still running at capture; this archive does not establish its present status. Earlier external Work experiments and the launch handoff remain historical, separate records.

## Preparing a genuinely shorter-path experiment

The [next native experiment task and offline acceptance tools](research/native-compression-next-2026-11-16/README.md) target a strict decrease in fully expanded primitive moves against frozen archived baselines. The two ordinary AC examples already have exact shortest-distance certificates; the Stable AC paths of 10 and 15 moves remain upper bounds. The literal macro expander and pinned official replay reject naming an unchanged path as a single macro, work-only improvements and altered baselines; 22 engineering tests pass. This is preparation for the original Yohaku team, not a new native launch, new shorter witness or platform submission. Live admission and pending-task checks must happen on the user's host.

## Independent round-two follow-up (2026-11-17)

The [new external review](research/native-round2-external-review-2026-11-17/README.md) independently rebuilds the complete Stable sac-08491 forward B4 and reverse B5 balls, compares every stored state, distance and witness, and verifies the existing ten-move upper bound. Together with the rank bound, this establishes distance exactly ten for the fixed instance. It does not turn the original native lead's error/interrupted task into a completed task. The pair-AND theory proofs pass a separate review and small independent checks; sac-07445 remains a fifteen-move upper bound. No shorter path, platform result, paired macro advantage or new native launch is claimed.

## Concept transformations and structural solving: research addendum

The [2026-11-17 user-directed supplement](research/concept-transformations-2026-11-17/README.md) asks the original native team to study equivalent representations, information retained for a specified solving goal, auxiliary information needed for recovery, and recognizable structures that yield solutions. It extends the earlier bounded concept-cost models; examples motivate the work without becoming new findings or a detour into unrelated set cardinalities. The [local Codex task](research/concept-transformations-2026-11-17/CODEX_TASK.md) preserves both research directions and historical records. This is a prepared handoff, not a new native dispatch or scientific result.

## Concept-transformation follow-up review (2026-11-17)

The [external follow-up review](research/native-concept-transformations-review-2026-11-17/README.md) checks the completed native supplement and independently reviews its Pair-AND inverse-image and quotient-walk lifting proofs. New small implementations pass for pair graphs through four vertices and quotient graphs through three; the full native enumerations and application beam search were not independently rerun. The application records no improvement over fifteen moves and no pruning by its matrix lower bound. This separates target-relevant information, executable path reconstruction and demonstrated computational benefit. Original model errors, stops and recovery records remain unchanged; no new native tasks were dispatched by the review.
````

## 输出与关联
```yaml
output: "README.md"
```

## AI 笔记
```yaml
{"source_version": 1, "writer": "Yohaku release packaging", "summary": "Portable release source. Scope and packaging adjustments are documented in PROVENANCE.md."}
```
